\chapter{Introduction}
\label{ch:introduction}

\section{Background and Context}
\label{sec:background}

Institutional kitchens such as college messes, hostel canteens, bakeries and
banquet halls routinely finish a service with edible food left over. At the same
time, shelters and community kitchens nearby are short of meals. Food waste is a
recognised problem at every scale~\cite{unepFoodWaste2024}. For a campus or a
tier-2/3 city, however, the binding constraint is rarely supply. It is
\emph{coordination under a deadline}: cooked food is safe for only a few hours,
and informal phone calls often outlast that window.

FoodLink is a web platform that brings every party to a handover into one tracked
workflow. A \textbf{donor} posts surplus food. An administrator verifies each
\textbf{recipient organisation} (an NGO or community kitchen) before it can
receive food. A \textbf{volunteer courier} collects and delivers it, and a
\textbf{platform administrator} maintains the platform. The project is developed
for UCS503P (2026--27 odd semester), which requires a proposal in week~4, this
prototype-stage (mid-term) report in week~7 and a final report in
week~17~\cite{ucs503Project}. The report describes what has been built against
the proposal~\cite{foodlinkProposal}, how it was verified, and what remains.

Although the repository is named \emph{FoodBridge-AI}, FoodLink contains
\textbf{no machine learning or computer vision}. Organisations are ranked by a
transparent weighted sum of five published criteria that can be recomputed by
hand. This choice is recorded as decision D-05 in the project's decision
log~\cite{foodlinkArchitecture}. All technical statements were checked against
the source code at commit \texttt{e588b35}~\cite{foodlinkRepo}.

\subsection{Problem Statement}
\label{sec:problem}

Informal coordination causes five problems. The first is
\textbf{perishability}: the hard deadline is routinely missed. The second is
\textbf{fragmentation}: donors call whoever they know, so availability is
invisible to others. The third is \textbf{lack of evidence}: nobody knows how
long claiming took, whether food arrived, or how much was lost. The fourth is
\textbf{the last mile}: courier availability is untracked. The fifth is
\textbf{trust and privacy}: donors need assurance that a real organisation is
receiving the food, and a home location should not be broadcast to anyone who
signs up. The problem is therefore to \emph{build a web platform that lets a
donor post surplus food with a hard deadline, directs it to nearby verified
organisations through an explainable ranking, coordinates a courier through
pickup and delivery, and records every step as server-stamped evidence, while
each participant sees and does only what its role and relationship to the
donation permit.}

\subsection{Motivation}
\label{sec:motivation}

Coordination, not supply, is the bottleneck. Shortening the delay between
posting and acceptance converts waste into meals, which is why the proposal's
primary metric is \emph{time-to-claim}~\cite{foodlinkProposal}. Trust must be
earned: anyone can register as an organisation, so FoodLink separates holding a
role from being trusted, and administrator verification gates who may receive
food. Decisions should be explainable, and a published weighted sum lets
organisations, donors and evaluators see \emph{why} a ranking came out as it
did, which an opaque model trained without outcome data would not. Finally,
evidence should not be self-reported. Metrics are credible only if the server
stamps their timestamps, so every transition is appended to a ledger.

\subsection{Existing System}
\label{sec:existing}

In the proposal's setting, redistribution relies on personal calls, messaging
groups and ad-hoc arrangements, and no shared software system
exists~\cite{foodlinkProposal}. Surplus is therefore visible only to the
contacts a donor happens to call. The recipient is whoever answers first rather
than the best-placed organisation. Deadlines are mentioned but not enforced,
and trust rests on familiarity rather than vetting. Couriers are arranged
separately and not tracked, nothing is recorded after handover, and addresses
are shared freely. The comparison is with informal practice in general, not
with specific third-party products.

\subsection{Proposed System}
\label{sec:proposed}

FoodLink is a three-tier web application: a React client, a REST API in Python
(FastAPI), and a relational database (SQLite by default). The API is the only
authority for every rule. A donation moves through the following flow:
\begin{enumerate*}[label=(\arabic*)]
  \item a donor posts food details, a pickup pin and an absolute deadline;
  \item the server stores the donation as \code{AVAILABLE}, ranks eligible
    organisations and, if any exist, records \code{MATCHED}, which is only a
    suggestion;
  \item a verified organisation, shown its own match explanation, accepts it
    (\code{ACCEPTED});
  \item a courier, who until then sees only a coarse area of about 1\,km, claims
    the pickup (\code{VOLUNTEER_ASSIGNED}) and only then sees the exact pin;
  \item the courier records \code{PICKED_UP} and \code{DELIVERED};
  \item the organisation confirms receipt (\code{COMPLETED}), which updates the
    reliability used in future rankings.
\end{enumerate*}
A donor (before pickup) or an administrator may also cancel. The organisation
may release a claimed pickup, and an administrator's expiry sweep closes
overdue, unaccepted donations. Every transition appends a server-stamped event.

\subsection{Scope}
\label{sec:scope}

\textbf{Implemented:} registration and sign-in, with the first administrator
created from the command line; role, ownership, lifecycle and verification
checks; donation posting with photo resizing; explainable ranking; the
nine-state lifecycle with race-safe writes and an event ledger; standing
requirements and a donor needs board; verification, the expiry sweep and
metrics; four web portals; rate limiting; and automated tests in continuous
integration.

\textbf{Partially implemented:} cancellation, pickup release, account
management and setting an organisation's map coordinates work only through the
API. Phone-layout screens at \code{/m/*} can be reached only by typing their
address.

\textbf{Not built:} notifications, maps and routing; object storage for photos;
scheduling of the expiry sweep; any deployment or real-user pilot; machine
learning or computer vision; and certification of food safety. The platform
records what a recipient needs in order to decide, but does not vouch for the
food~\cite{foodlinkProposal}.

\section{Project Objectives}
\label{sec:objectives}

Each objective states a measurable criterion and its verified mid-term status.
\begin{enumerate}
  \item \textbf{Objective 1:} Donors post surplus food with an absolute,
    server-validated deadline. \emph{Met.}
  \item \textbf{Objective 2:} Eligible organisations are ranked by a published
    weighted heuristic that returns per-criterion scores and reasons.
    \emph{Met.}
  \item \textbf{Objective 3:} The lifecycle refuses illegal and conflicting
    transitions and records one event per transition. \emph{Met; covered by
    lifecycle and concurrency tests.}
  \item \textbf{Objective 4:} Every read and write is restricted by role,
    ownership, state and verification, and sensitive data is withheld from
    readers who do not need it. \emph{Met; covered by scope tests.}
  \item \textbf{Objective 5:} Organisations publish standing requirements that
    donors can browse. \emph{Met.}
  \item \textbf{Objective 6:} Time-to-claim, rescue rate, expiry-loss rate and
    handover time are derived from the event ledger. \emph{Instrumented; not yet
    measured on real use.}
  \item \textbf{Objective 7:} Automated tests run in CI on every push.
    \emph{Met: 495 backend and 156 frontend tests pass.}
  \item \textbf{Objective 8:} A staging deployment and pilot before week~17.
    \emph{Not started.}
\end{enumerate}

\section{Organisation of the Report}
\label{sec:organisation}

The report follows the TIET report template~\cite{tietTemplate}.
Chapter~\ref{ch:methodology} covers the methodology, architecture, requirements,
design and implementation modules. Chapter~\ref{ch:results} presents testing,
results and limitations, and Chapter~\ref{ch:feasibility} is the Feasibility
Report. Chapter~\ref{ch:conclusion} concludes, and Chapter~\ref{ch:uml} contains
the use case diagram and scenarios, the class diagram and the activity diagram.
